\documentclass{article}
\usepackage{cite}

\begin{document}

\title{A Survey on Mobile Edge Computing}
\author{Test Author}
\maketitle

\section{Introduction}

Mobile edge computing (MEC) has emerged as a promising paradigm to address the increasing demand for low-latency and computation-intensive applications \cite{qian2022game}. Recent studies have shown that deep reinforcement learning can be effectively applied to optimize resource allocation in connected vehicles \cite{he2017integrated}.

The integration of edge computing with vehicle networks has become a hot research topic in recent years. Various approaches have been proposed to tackle the challenges of task offloading, resource allocation, and latency optimization.

\section{Related Work}

\subsection{Task Offloading in Edge Computing}

Various game-theoretic approaches have been proposed for task offloading in MEC environments. Qian et al. proposed a game theory based D2D collaborative offloading scheme \cite{qian2022game}, which demonstrated significant improvements in workflow execution time. Their approach considers the cooperative behavior between devices and formulates the offloading problem as a non-cooperative game.

\subsection{Deep Learning for Connected Vehicles}

The integration of networking, caching, and computing has been studied in the context of connected vehicles. He et al. \cite{he2017integrated} proposed a deep reinforcement learning approach that jointly optimizes these three aspects. Their work shows that DRL can effectively handle the dynamic nature of vehicular networks.

\subsection{Other Approaches}

Some researchers have explored alternative methods. According to some claims \cite{fake2023paper}, there exist methods that can achieve 100\% efficiency, but these results have not been verified.

\section{Methodology}

Our approach builds upon the insights from both game-theoretic \cite{qian2022game} and deep learning methods \cite{he2017integrated} to propose a hybrid solution for mobile edge computing optimization.

\section{Conclusion}

In this paper, we have surveyed the recent advances in mobile edge computing and connected vehicles.

\bibliographystyle{IEEEtran}
\bibliography{sample}

\end{document}
